\documentclass[UTF8, 11pt]{ctexart}
\usepackage[a4paper, margin=1in]{geometry}
\usepackage{amsmath, amssymb, bm}
\usepackage{hyperref}

\title{Stage-5 部署数学:\\
       延迟预算、成功率转换、Stateless vs Stateful 选择}
\author{REACT stage-5 启动前数学}
\date{2026-05-22}

\newcommand{\R}{\mathbb{R}}
\newcommand{\ms}{\,\text{ms}}
\newcommand{\flops}{\,\text{FLOPS}}

\begin{document}
\maketitle

英文版:\href{run:./04_stage5_deployment_math.tex}{04\_stage5\_deployment\_math.tex}

\section{论文要的两个目标数字}

REACT 论文级目标(README 里设的,源自 PEMTRS RA-L 2026):
\begin{enumerate}
  \item 推理延迟 $T_{\text{inf}} < 10\ms$ / 帧,在 Jetson 级设备
        (Orin Nano 或 NX 是典型参考)。
  \item 闭环成功率 $\text{SR} \geq 85\%$,在动态障碍场景。
\end{enumerate}

当前训练指标 \texttt{dyn\_dyn\_tail} $\approx 0.22$--$0.24$ **既不是这两个**。
本文档推导这个桥。

\section{延迟预算分解}

每帧推理,在 temporal 模式(K=10 stateless)下:
\begin{equation}\label{eq:latency-stateless}
T_{\text{inf}}^{\text{stateless}}
\;=\;
T_{\text{ReVAE-enc}}^{K} + T_{\text{GRU}}^{K} + T_{\text{backbone}}^{1} +
T_{\text{head}}^{1} + T_{\text{post}}.
\end{equation}

单帧模式(stage-3.1):
\begin{equation}\label{eq:latency-single}
T_{\text{inf}}^{\text{single}}
\;=\;
T_{\text{ReVAE-enc}}^{1} + T_{\text{backbone}}^{1} +
T_{\text{head}}^{1} + T_{\text{post}}.
\end{equation}

Stateful 模式(DiffPhysDrone 风格,见
\texttt{stage\_3\_4\_research\_cn.md}):
\begin{equation}\label{eq:latency-stateful}
T_{\text{inf}}^{\text{stateful}}
\;=\;
T_{\text{ReVAE-enc}}^{1} + T_{\text{GRU}}^{1} + T_{\text{backbone}}^{1} +
T_{\text{head}}^{1} + T_{\text{post}}.
\end{equation}

\subsection{各组件估算(Orin NX,INT8 量化的 reVAE)}

按
\href{https://developer.nvidia.com/embedded/jetson-orin}{Jetson Orin NX} INT8
吞吐率(FP16 baseline $\approx 70$ TFLOPS)和模块 FLOP 数(\texttt{reVAE.encode}
在 $96\times160$ 输入下 $\approx 0.4$ GFLOPS / 帧;\texttt{YopoBackbone}
ResNet18 在同样输入下 $\approx 1.8$ GFLOPS):
\begin{center}
\begin{tabular}{lrrr}
\hline
组件 & FLOPS / 次 & 时间 / 次 (FP16) & 时间 / 次 (INT8) \\
\hline
\texttt{reVAE.encode}     & 0.4 G & 0.50\ms & 0.20\ms \\
\texttt{reVAE.decode}     & 0.4 G & 不在推理图上 & --- \\
\texttt{nn.GRU} 一步      & 0.1 M & 0.05\ms & 0.05\ms \\
\texttt{YopoBackbone}     & 1.8 G & 2.4\ms  & 1.0\ms \\
\texttt{YopoHead}         & 5.0 M & 0.08\ms & 0.04\ms \\
后处理(poly solve、anchor argmax) & --- & 1.0\ms & 1.0\ms \\
\hline
\end{tabular}
\end{center}

\subsection{延迟预算表}

代入(FP16):
\begin{center}
\begin{tabular}{lrl}
\hline
模式                    & 延迟 & 打到 $<10\ms$ 预算? \\
\hline
\texttt{single}(3.1)   & $0.5 + 2.4 + 0.08 + 1.0 \approx 4.0\ms$  & **是**(还有余量) \\
\texttt{stateful} K=1   & $0.5 + 0.05 + 2.4 + 0.08 + 1.0 \approx 4.0\ms$ & **是** \\
\texttt{stateless} K=10 & $5.0 + 0.5 + 2.4 + 0.08 + 1.0 \approx 9.0\ms$  & **勉强**(10\% 余量) \\
\hline
\end{tabular}
\end{center}

结论:**stage-3.4 stateless K=10 在 Orin NX FP16 路径上可行但不舒服**。
要么走 INT8(reVAE encode 每次降到 0.2\,ms,K=10 总时 $\sim 6.0\ms$),
要么部署时换 stateful(保持 $\sim 4.0\ms$)。Stateful 重构**不是**模型
重训 —— 它是推理循环的改动:逐帧喂 depth,GRU hidden 在调用之间保持。

\section{从 \texttt{dyn\_dyn\_tail} 到成功率}

这是欠定的桥。我们有
\[
\texttt{dyn\_dyn\_tail} = 0.23 \quad \text{(stage-3.1+,v3 数据集)}.
\]
我们想预测闭环 sim 里的 $\text{SR}$。没有干净的闭式表达,因为 $\text{SR}$
依赖于 (i) 闭环策略稳定性,(ii) 动态场景难度,(iii) 控制器跟踪误差。
但我们可以**给出界**。

\subsection{用损失率先验估计每轨碰撞概率}

记 $p_c$ 为每帧预测 endstate $\bm{p}_i^{\text{pred}}$ 落进禁止区(离某个障碍
$d_{\text{safe}}$ 之内)的概率。假设最优 anchor 以高概率被选中,而
\texttt{dyn\_dyn} 大致是激活 anchor 上罚款的均值:
\begin{equation}\label{eq:pc-estimate}
\texttt{dyn\_dyn}
\;\approx\;
\lambda_{\text{dyn}} \cdot p_c \cdot \overline{\softplus} ,
\end{equation}
其中 $\overline{\softplus}$ 是激活 anchor 上的 softplus 均值 $\approx 0.3$。
$\lambda_{\text{dyn}} = 3.0$、\texttt{dyn\_dyn} $= 0.23$ 给出
\begin{equation}
p_c \approx \frac{0.23}{3.0 \cdot 0.3} \approx 0.26.
\end{equation}
解释:$\approx 26\%$ 的帧里,**某个**候选 anchor 会出违反安全边界的 endstate。
但被选中的(top-score)anchor 通常是更安全的,所以每帧**实际**碰撞概率
显著更低。

\subsection{转换到闭环 $\text{SR}$}

假设一个闭环轨迹由 $N_{\text{frame}} \approx 100$ 个规划步组成
(30Hz 下 3.3\,s 飞行),score head 每次选**最优** anchor(所以每步**实际**
碰撞概率 $p_{\text{actual}} \ll p_c$,可能小一个数量级,因为 score 选择)。
那么
\begin{equation}\label{eq:sr-bound}
\text{SR}
\;\approx\;
(1 - p_{\text{actual}})^{N_{\text{frame}}}
\;=\;
\exp\!\bigl(-N_{\text{frame}} \cdot p_{\text{actual}}\bigr).
\end{equation}
要达到 $\text{SR} = 0.85$,需要
$p_{\text{actual}} \le -\ln(0.85)/100 \approx 1.63 \times 10^{-3}$。
$\text{SR} = 0.90$ 时:$p_{\text{actual}} \le 1.05 \times 10^{-3}$。

如果 $p_{\text{actual}} \approx p_c/k$,其中 $k$ 是 score-selection 改善因子,
那需要 $k \ge p_c / (1.63\times 10^{-3}) = 0.26 / 1.63\times 10^{-3}
\approx 160$。

\subsection{Score head 改善因子 $k$}

**这个值不跑闭环测不出来**。经验上 YOPO score head 在静态场景下 $k$
量级在 $50$--$200$(baseline 报告 $\sim$$90\%$ 静态成功率);是否能延续到
动态场景**正是 stage-5 要测的**。

**所以"\texttt{dyn\_dyn} = 0.23 够不够 $\text{SR} \geq 85\%$"的答案是:
取决于 $k$**,**不跑 stage-5 测不出来**。

\section{决定:Stage-5 先,不会 block 监督改进}

用户的问题:"如果我直接走 stage-5,会不会 block $<10\ms$ 和 $\geq 85\%$ 目标?"

**不会 block;它恰恰是测这两个目标的唯一方法**。理由:

\begin{enumerate}
  \item **$<10\ms$ 目标**依赖推理图 + 量化 + 硬件选择(Jetson NX vs Orin Nano)。
        Stage-5 **就是**这个工作。监督改动不会移动这个数字;Option B head
        $N=5$ 最多给 head 加 $\sim$$0.2\ms$(可忽略)。Stage-5 先决定
        stateless K=10 在预算内还是要换 INT8 / stateful。**这两个都不需要重训**。
  \item **$\geq 85\%$ SR 目标**需要现在不存在的闭环 sim 基础设施。建这个
        是 stage-5 的一部分。建好之后跑当前 stage-3.1 模型,测 $k$ 和实际
        $\text{SR}$。**如果**测出来 $\text{SR} \geq 85\%$,**完全不需要**改
        监督 —— ship + 写论文。**如果**测出来 $\text{SR} < 85\%$,这个 gap
        告诉你监督改动该多激进(多 waypoint vs ESDF vs 数据 vs 都改)。
  \item **没有闭环 bench 的监督升级是投机**。把 \texttt{dyn\_dyn} 从 $0.23$
        通过 Option B 拉到 $0.16$,只有在 $\text{SR}$ 真的低于 $85\%$
        **且**瓶颈是规划器的避障质量(不是控制器,不是 FOV 覆盖问题)时
        才值。**没测**就动手,可能两周做 Option B 之后发现 $k$ 在
        $\texttt{dyn\_dyn} = 0.23$ 下已经够用了,白干。
\end{enumerate}

\subsection{推荐顺序}

\begin{enumerate}
  \item **Stage-5.A**:当前 stage-3.1 模型的 ONNX/TensorRT 导出 + Orin NX
        延迟 profile。约 3 天。
  \item **Stage-5.B**:用 \texttt{Simulator/} 已有的 ROS 控制器写闭环 sim
        harness。跑 100 个动态场景,测 $\text{SR}$。约 4 天。
  \item **决策点**。$\text{SR} \geq 85\%$:写论文 + ship。$\text{SR} < 85\%$:
        按 \texttt{01\_collision\_loss\_saturation\_cn.tex} 的分析挑期望收益
        最大的监督升级(很可能 Option B,可能加 camera-aware bake)。
\end{enumerate}

这个顺序**信息每周产出最大**:stage-5 单独要么直接给出 paper,要么明确告诉你
该追哪个监督改动。先做监督改动是反过来 —— polish 一个跟 paper-target SR
转换不确定的指标。

\end{document}
